\begin{musica}{Porto, Pátria, Portugal}{Tuna de Engenharia da Universidade do Porto}{}
\nota{(vocal)} Porto, Pátria, Portugal

\nota{Refrão:}
\ac{D}{}_ Meu \ac{G}{Por}to velhinho, encarquilha\ac{E}{di}nho,
Dos tempos da Funda\ac{Am}{ção};
Ressurges, o\ac{D7}{van}te,
Belo e triunfante nas festas de S. Jo\ac{G}{ão}.
\ac{D7}{Dan}ças e fol\ac{G}{gue}dos, amor, sonhos \ac{E}{le}dos,
E cascatas, um en\ac{Am}{can}to;
Foguetes no \ac{G}{ar}, lindos a esta\ac{D7}{lar}
Em louvor do meigo \ac{G}{San}to.

\ac{Gm}{}__ Portugal nasceu de um sonho
Famoso, heróico e risonho,
Do seu Régio Funda\ac{D7}{dor};
Que, a golpes de montante,
Do férreo Mouro arrogante,
Foi duro Conquista\ac{Gm D}{dor}. _
\ac{Gm}{}__ Nas suas lutas brilhantes,
Sempre duras, triunfantes,
A Virgem tomou por \ac{Cm}{guia};
E às terras que conquis\ac{Gm}{tou},
Logo com amor cha\ac{D7}{mou},
“Terras de Santa Ma\ac{G}{ria}”.

Tornou, do “Porto e Cale”,
O nome que muito vale,
Desta Pátria de eleição;
E, p’ra mais nobre a tornar,
Deu-lhe - a riqueza sem par -
O Porto por coração!

Pátria, Porto, Portugal!
São três nomes sem igual,
Nomes de rara beleza;
Portugal, lindo recanto,
Pátria, mãe, és encanto;
Porto és honra e nobreza!

\posicoes{Cm: x35543, D: xx0232, D7: xx0212, E: 022100, G: 320033, Gm: 355333,
  Am: x02210}
\end{musica}
